\section*{Part IV: Leasing versus Owning: The Actual Economics}
\addcontentsline{toc}{section}{Part IV: Leasing versus Owning: The Actual Economics}


\subsection{29. Capital cost, operating cost and option value}

Compare owning and leasing on a \textbf{like-for-like present-value basis}. Include deposits, stamp duty or lease transaction costs, rent escalation, fit-out, repairs, insurance, property taxes, facilities management, financing, residual value, and restoration costs. Ownership benefits from potential appreciation and retained improvements but requires more capital and management. Lease economics depend on renewal and landlord risk.

Never compare 1st-year rent with land price alone. Property ownership and campus operations are different businesses with different risk profiles.


\subsection{30. A 5-year leased studio example}

\textbf{[Model]} A 2,000-square-foot leased studio has base rent of \textnormal{₹}1.2 lakh/month, increasing 5% annually; an upfront refundable security deposit of \textnormal{₹}6 lakh; fit-out cash expense \textnormal{₹}12 lakh; and annual maintenance/utilities of \textnormal{₹}5 lakh escalating 5%. Ignoring taxes and discounting, 5-year nominal rent totals approximately \textbf{\textnormal{₹}79.57 lakh}; maintenance totals \textbf{\textnormal{₹}27.63 lakh}. Together with fit-out, the 5-year nominal expense is about \textbf{\textnormal{₹}119.20 lakh}, excluding deposit financing cost, restoration and other overhead. The deposit is a cash commitment, not automatically an expense if recoverable.

This is not a Mumbai or Pune market quotation. Replace every input with binding local quotes. Calculate NPV using a cost of capital appropriate to the business and a credible lease renewal scenario.


\subsection{31. An owned-space example}

\textbf{[Model]} Suppose comparable premises cost \textnormal{₹}2.4 crore to acquire and fit out; annual operating and maintenance cost is \textnormal{₹}6 lakh; the asset could be sold after 5 years at \textnormal{₹}2.4 crore before transaction costs. At a 12% discount rate, a simplified unlevered present-value occupancy cost is initial \textnormal{₹}240 lakh plus PV of \textnormal{₹}6 lakh annual expenses minus PV of \textnormal{₹}240 lakh terminal proceeds: approximately \textnormal{₹}240 + 21.63 - 136.18 = \textbf{\textnormal{₹}125.45 lakh}. This ignores taxes, legal costs, appreciation uncertainty, rental alternatives, capital repairs and sale costs.

The numerical closeness to the lease example should not be interpreted as a recommendation. Change the terminal sale price, utilization, financing, and tax assumptions: the answer can reverse. Unlike a lease, ownership ties up capital that might earn more in faculty, sales, and research.


\subsection{32. Why the hurdle rate differs for a campus}

A founder may view real estate as ``safe'' because it is tangible. For a business, the relevant cost includes liquidity risk, debt obligations, concentration, alternative investments, and the danger of having to sell during a downturn. A nominal 8% property appreciation assumption can be dominated by missed operating opportunities and forced-sale discounts.

Use comparable risk-adjusted NPV with scenarios. Treat any nonfinancial mission benefit separately and explicitly rather than hiding it inside an optimistic sale value.


\subsection{33. A lease negotiation checklist}

Negotiate permitted educational/research use, fit-out rights, lock-in period, rent escalation, deposit safeguards, maintenance allocation, utility capacity, accessibility, signage, subleasing, assignment, early termination, renewal, restoration, repair access, insurance, property-title warranties, operating hours and fire-safety responsibilities. A laboratory lease may need waste, ventilation, power load and hazardous-material clauses.

A nominally affordable long lease can be more irreversible than a small property purchase if termination penalties and fit-out losses are severe.


\subsection{34. The build-to-suit alternative}

A landowner or developer may construct premises to a specification in exchange for a long lease or service charge. This can preserve upfront cash and offload some development responsibility. The tradeoff is contractual rigidity, financing costs embedded in rent, insolvency/step-in exposure, and possible delays. Negotiate completion acceptance tests, change-control rules, defects liability, maintenance obligations, and remedies for late delivery.

Only select build-to-suit when the institute can credibly support the lease even if enrollment or grants are below plan.
